\begin{table}[htbp]
\centering
\caption{Summary Statistics by Treatment Group}
\label{tab:summary}
\begin{adjustbox}{max width=\textwidth}
\begin{threeparttable}
\begin{tabular}{lcccc}
\toprule
 & \multicolumn{2}{c}{Participating States} & \multicolumn{2}{c}{Non-Participating (MA, CT, IL)} \\
\cmidrule(lr){2-3} \cmidrule(lr){4-5}
 & Pre-Treatment & Exposed & Pre-Treatment & Exposed \\
 & (Born 1912--1921) & (Born 1922--1928) & (Born 1912--1921) & (Born 1922--1928) \\
\midrule
N & 6,300,443 & 5,179,082 & 823,982 & 656,117 \\
Mean Wage Income (\$1950) & 2532 & 2088 & 2651 & 2152 \\
SD Wage Income & 1612 & 1442 & 1584 & 1385 \\
Mean Education (years) & 7.7 & 8.5 & 8.4 & 9.3 \\
Mean Occ. Income Score & 26.0 & 24.2 & 27.5 & 25.6 \\
Pct Male & 73.1 & 71.2 & 71.0 & 67.9 \\
Pct White & 94.6 & 92.8 & 99.2 & 98.3 \\
Pct Rural (1930) & 32.5 & 29.2 & 12.8 & 10.5 \\
\bottomrule
\end{tabular}
\begin{tablenotes}[flushleft]\footnotesize
\item \textit{Notes:} Data from the IPUMS Multigenerational Longitudinal Panel (MLP), linking individuals across the 1930, 1940, and 1950 full-count censuses. Participating states accepted federal Sheppard-Towner funds (1921--1929). Non-participating states (MA, CT, IL) refused on states' rights grounds. Wage income and occupational income score measured in 1950.
\end{tablenotes}
\end{threeparttable}
\end{adjustbox}
\end{table}
